\chapter{Bahasa: dari Kantong Kata ke Model Bahasa Besar}
\label{ch:nlp}

Kuliah natural language processing saya ikuti pada musim dingin 2025/26 di Institute of Computational Perception. Sepuluh
pertemuannya bergerak dari pra-pemrosesan teks dan representasi dokumen, representasi kata dan embedding, jaringan saraf,
model bahasa $n$-gram, embedding kata dan kalimat yang dipelajari, attention dan transformer, NLP multimodal, sampai model bahasa
besar. Tiga tugas kelompoknya mengikuti alur yang sama: mengklasifikasikan dokumen dengan metode machine learning standar,
menjelajahi word embedding, lalu mengklasifikasikan cerita pendek ke banyak label sekaligus dengan PyTorch dan BERT.

Transformer dan model bahasa besar sudah muncul di \cref{ch:lstm} dan \cref{ch:terbaru}. Bab ini mengisi jalan panjang di
antaranya: bagaimana teks diubah menjadi angka, dan bagaimana gagasan-gagasan sederhana tentang menghitung kata perlahan tumbuh
menjadi model yang kini dipakai jutaan orang.

\section{Membersihkan teks}

Sebelum apa pun, teks harus diseragamkan. Huruf besar dikecilkan, sehingga ``Linz'' dan ``linz'' menjadi satu kata, walau
kadang dengan harga ambiguitas, seperti \emph{General Motors} dan \emph{general motors}. Angka dan tanggal sering diganti token
khusus, karena kalau setiap angka menjadi kata tersendiri, kosakata meledak. Kata majemuk dalam bahasa seperti Jerman perlu
dipecah, \emph{Halsschlagader} menjadi leher, pukulan, dan pembuluh. Kata-kata yang umum, \istilah{stop words} seperti
``yang'', ``dan'', ``di'', sering dibuang untuk pendekatan kantong kata, walau dalam frasa seperti ``to be or not to be'' justru
merekalah isinya.

Variasi bentuk kata dikurangi dengan dua cara. \istilah{Stemming} memotong imbuhan dengan aturan kasar. Algoritme Porter
\citep{porter1980}, yang paling umum untuk bahasa Inggris, memakai aturan seperti \emph{sses} menjadi \emph{ss}, \emph{ies}
menjadi \emph{i}, dan \emph{ational} menjadi \emph{ate}, sehingga \emph{caresses} menjadi \emph{caress}, \emph{ponies} menjadi
\emph{poni}, dan \emph{relational} menjadi \emph{relate}. Hasilnya tidak harus kata yang sungguhan. \istilah{Lemmatization} memakai
kamus untuk mengembalikan bentuk dasar yang benar, misalnya \emph{am}, \emph{are}, dan \emph{is} menjadi \emph{be}. Lebih teliti,
tetapi lebih lambat, dan tidak berdaya untuk kata yang tidak ada di kamus. Bahasa Indonesia, dengan imbuhan seperti
\emph{me-}, \emph{-kan}, dan \emph{per-an}, menghadapi masalah yang sama dalam bentuk yang lebih kaya.

\section{Kantong kata}

Representasi paling sederhana untuk sebuah dokumen adalah \istilah{bag of words}: vektor sepanjang kosakata yang berisi berapa
kali setiap kata muncul, dengan urutan dibuang sama sekali. ``Anjing menggigit orang'' dan ``orang menggigit anjing'' menjadi
vektor yang sama. Meski kasar, representasi ini, ditambah pembobotan TF-IDF dari \cref{ch:rekomendasi}, sudah cukup untuk banyak
tugas klasifikasi dokumen. Tugas pertama kuliah membangun pengklasifikasi sentimen seperti ini dengan regresi logistik dan
pengklasifikasi standar dari \cref{ch:unsup}.

Masalah utamanya adalah \emph{sparsity}. Kosakata berisi puluhan ribu kata, sedangkan satu kalimat hanya memakai belasan, sehingga
vektornya hampir seluruhnya nol. Lebih parah lagi, setiap kata adalah dimensi yang berdiri sendiri. ``Bagus'' dan ``baik''
sama jauhnya dengan ``bagus'' dan ``buruk''. Model tidak bisa memindahkan apa yang ia pelajari tentang satu kata ke kata yang
bermakna serupa.

\section{Makna dari tetangga}

Jalan keluarnya datang dari linguistik. Firth menulis pada 1957 bahwa kita mengenal sebuah kata dari teman-temannya. Kuliah
memakai contoh kata yang hampir tidak dikenal, \emph{tesgüino}. Dari kalimat-kalimat di sekitarnya, bahwa ia diminum, difermentasi
dari jagung, dan dibuat untuk upacara, kita menyimpulkan bahwa ia sejenis minuman beralkohol, tanpa pernah melihat definisinya.
Dua kata bermakna mirip kalau mereka sering muncul di tengah kata-kata yang sama. Inilah \istilah{distributional semantics}.

Gagasan ini dihitung dengan matriks kata--konteks: baris untuk setiap kata, kolom untuk setiap kata konteks, dan isinya berapa
kali keduanya muncul berdekatan, misalnya dalam jendela beberapa kata. Hitungan mentah didominasi kata-kata umum, sehingga
diganti dengan \istilah{pointwise mutual information},
\begin{equation}
\mathrm{PMI}(w,c)=\log_2\frac{p(w,c)}{p(w)\,p(c)},
\end{equation}
yang bertanya seberapa sering dua kata muncul bersama dibanding yang diharapkan kalau keduanya independen. Misalkan dalam seribu
pasangan kata--konteks, pasangan ``kopi''--``seduh'' muncul 10 kali, ``kopi'' muncul 50 kali sebagai kata, dan ``seduh'' 40 kali
sebagai konteks. Peluang bersamanya $0{,}01$, sedangkan hasil kali peluang terpisahnya $0{,}05\times0{,}04=0{,}002$. PMI-nya
$\log_25\approx2{,}3$: kedua kata muncul bersama lima kali lebih sering dari kebetulan. Nilai negatif sulit ditaksir dengan andal
dari data terbatas, sehingga biasanya dipotong ke nol, \istilah{positive PMI}. Kuliah juga mencatat kelemahan PMI: ia terlalu
menghargai kata yang jarang, karena satu kebetulan saja bisa memberi rasio yang besar.

Matriks PPMI lebar sekali dan jarang. Memampatkannya dengan SVD (\cref{lamp:matematika}) dan mengambil beberapa ratus dimensi
pertama memberi vektor padat untuk setiap kata, sebuah \istilah{word embedding}. GloVe \citep{pennington2014} merumuskan hal serupa
sebagai masalah regresi: cari vektor kata $\bm w_i$ dan vektor konteks $\tilde{\bm w}_j$ beserta bias sehingga
\begin{equation}
\bm w_i^\top\tilde{\bm w}_j+b_i+\tilde b_j\approx\log X_{ij},
\end{equation}
dengan $X_{ij}$ hitungan kemunculan bersama, diminimalkan dengan galat kuadrat berbobot. Bobotnya meredam pasangan yang sangat
sering dan mengabaikan pasangan yang tidak pernah muncul. Hasil kali dalam yang mendekati logaritma hitungan bersama, setelah bias
memperhitungkan frekuensi masing-masing kata, pada dasarnya adalah PMI. Dari sinilah sifat analogi yang terkenal muncul: vektor
dari ``pria'' ke ``wanita'' kira-kira sejajar dengan vektor dari ``raja'' ke ``ratu''. Tugas kedua kuliah menjelajahi sifat-sifat ini:
kemiripan kosinus, tetangga terdekat, dan analogi.

\section{Model bahasa}

Model bahasa memberi peluang pada deret kata. Dengan aturan rantai, peluang sebuah kalimat adalah hasil kali peluang setiap kata
diberikan kata-kata sebelumnya. Model $n$-gram menyederhanakannya dengan asumsi Markov: setiap kata hanya bergantung pada $n-1$ kata
sebelumnya, dan peluangnya ditaksir dari hitungan. Di \cref{ch:lstm} kita melihat contoh bigram ``nasi goreng'': kalau ``nasi''
muncul 200 kali dan ``nasi goreng'' 50 kali, peluangnya $0{,}25$.

Masalahnya adalah pasangan yang tidak pernah muncul di korpus, yang mendapat peluang nol, sehingga seluruh kalimat yang memuatnya
juga berpeluang nol. \istilah{Laplace smoothing} menambahkan satu ke setiap hitungan,
\begin{equation}
P_{\mathrm{Lap}}(w_t\mid w_{t-1})=\frac{\mathrm{hitung}(w_{t-1},w_t)+1}{\mathrm{hitung}(w_{t-1})+V},
\end{equation}
dengan $V$ ukuran kosakata. Untuk kosakata sepuluh ribu kata, peluang ``goreng'' setelah ``nasi'' turun dari $0{,}25$ menjadi
$51/10\,200=0{,}005$. Penghalusan ini terlalu kasar: hampir seluruh massa peluang dipindahkan ke sembilan ribu lebih kelanjutan
yang tidak pernah terlihat. Cara yang lebih baik adalah \istilah{backoff}, yang turun ke $n$-gram yang lebih pendek kalau hitungan
$n$-gram panjang tidak ada, atau \istilah{interpolation}, yang mencampur taksiran bigram dan unigram dengan bobot $\lambda$.

Kinerja model bahasa diukur dengan \istilah{perplexity}, kebalikan rata-rata geometris peluang per kata pada data uji,
$\mathrm{PP}=P(w_1,\dots,w_N)^{-1/N}$. Kalau model memberi peluang $0{,}1$, $0{,}2$, dan $0{,}05$ pada tiga kata sebuah kalimat, hasil
kalinya $0{,}001$, rata-rata geometrisnya $0{,}1$, dan perplexity-nya 10: model sama bingungnya dengan memilih seragam di antara sepuluh
kata. Model yang menebak seragam dari seluruh kosakata punya perplexity $V$. Perplexity adalah eksponen dari cross-entropy rata-rata,
sehingga menurunkannya sama dengan menurunkan loss pelatihan di \cref{ch:data}.

Model bahasa neural \citep{bengio2003} mengganti tabel hitungan dengan jaringan. Setiap kata konteks dipetakan ke embedding,
embedding-embedding itu digabung, dilewatkan ke lapisan tersembunyi, lalu ke softmax atas seluruh kosakata. Dengan kosakata sepuluh
ribu, embedding berdimensi 100, tiga kata konteks, dan 200 neuron tersembunyi, tabel embedding berisi satu juta parameter, lapisan
tersembunyi $300\times200=60\,000$, dan lapisan keluaran $200\times10\,000=2$ juta. Lapisan keluaran yang paling mahal, dan masalah
yang sama muncul lagi pada model besar. Keuntungannya, kata-kata yang mirip mendapat embedding yang mirip, sehingga apa yang
dipelajari dari ``kopi panas'' ikut membantu ``teh panas'', sesuatu yang tidak bisa dilakukan tabel hitungan.

Loss-nya adalah negative log likelihood setelah softmax, dan kuliah menurunkan gradiennya terhadap logit: $\bm p-\bm y$, peluang
yang diramalkan dikurangi label one-hot, bentuk yang sama dengan regresi logistik di \cref{ch:data}. Softmax menciptakan
persaingan di antara semua kata. Menaikkan logit satu kata otomatis menurunkan peluang semua kata lain.

\section{word2vec dan embedding kalimat}

word2vec \citep{mikolov2013} menyederhanakan model bahasa neural menjadi tugas yang hanya bertujuan menghasilkan embedding. Skip-gram
menebak kata-kata konteks dari satu kata pusat. Softmax atas seluruh kosakata terlalu mahal, sehingga dipakai
\istilah{negative sampling}: tugas klasifikasi banyak kelas diubah menjadi klasifikasi biner. Untuk pasangan sungguhan antara kata
pusat $c$ dan kata konteks $o$, serta $k$ kata acak $n_1,\dots,n_k$ yang ditarik dari distribusi derau,
\begin{equation}
\Loss=-\log\sigma\big(\bm u_o^\top\bm v_c\big)-\sum_{i=1}^k\log\sigma\big(-\bm u_{n_i}^\top\bm v_c\big).
\end{equation}
Suku pertama menaikkan skor pasangan sungguhan, suku kedua menurunkan skor pasangan palsu. Kalau skor pasangan sungguhan $2$, suku
pertamanya $-\ln\sigma(2)\approx0{,}13$. Kalau satu kata derau punya skor $1$, sumbangannya $-\ln\sigma(-1)\approx1{,}31$, hukuman besar
yang mendorong skor itu turun. Kata-kata derau biasanya ditarik dengan peluang sebanding dengan frekuensinya dipangkatkan $3/4$, agar
kata yang umum tidak terlalu mendominasi dan kata langka sesekali ikut terpilih.

Embedding kata belum tentu menjadi embedding kalimat yang baik. Merata-ratakan vektor kata dalam kalimat, cara paling sederhana,
membuang urutan dan memberi bobot yang sama pada kata ``yang'' dan kata ``gempa''. Pembobotan dengan TF-IDF memperbaikinya sedikit.
Kuliah membahas cara yang lebih baik, yaitu melatih model langsung dengan tujuan tingkat kalimat, misalnya meramalkan kata yang
dihilangkan dari kalimat memakai representasi seluruh kalimat lainnya, sehingga embedding yang dihasilkan memang dirancang untuk
mewakili kalimat utuh.

\section{Model bahasa besar}

Kuliah membagi model bahasa besar yang dilatih sebelumnya ke dalam tiga keluarga. Model encoder seperti BERT \citep{devlin2019}
melihat seluruh kalimat sekaligus dan menghasilkan embedding kontekstual: kata ``bank'' di ``bank sungai'' dan di ``bank sentral''
mendapat vektor yang berbeda. Model decoder seperti GPT \citep{brown2020} menulis dari kiri ke kanan. Model encoder--decoder
menggabungkan keduanya untuk tugas seperti terjemahan.

BERT dilatih dengan \istilah{masked language model}: sebagian token dalam kalimat ditutup, dan model menebak token yang ditutup dari
konteks di kedua sisinya. Contoh di kuliah: ``Jim made [MASK] for his girlfriend and [MASK] was very proud'', dengan jawaban
\emph{spaghetti} dan \emph{he}. Kosakatanya sekitar 30 ribu token, karena kata dipotong menjadi potongan-potongan dengan WordPiece,
saudara dekat byte-pair encoding di \cref{ch:terbaru}. BERT-Base punya 12 lapisan, dimensi tersembunyi 768, dan sekitar 110 juta
parameter. Slide kuliah memberi pembanding yang menarik: tabel embedding statis seperti word2vec dengan kosakata 200 ribu kata dan
vektor berdimensi 768 sudah berisi sekitar 153 juta parameter, lebih banyak dari seluruh BERT-Base. Tokenisasi subkata bukan hanya
mengatasi kata yang tidak dikenal, tetapi juga jauh lebih hemat.

Untuk menyetel BERT ke tugas klasifikasi, embedding keluaran dari token khusus \texttt{[CLS]} di awal kalimat dilewatkan ke lapisan
klasifikasi kecil, dan seluruh jaringan disetel ulang dengan laju belajar kecil. Untuk tugas pasangan kalimat, seperti menilai
apakah dua kalimat bermakna sama, keduanya digabung dengan token pemisah \texttt{[SEP]}. Tugas ketiga kuliah memakai resep ini untuk
mengklasifikasikan cerita pendek. Di situ muncul perbedaan antara \emph{multi-class} dan \emph{multi-label}. Kalau setiap cerita
punya tepat satu genre, keluaran memakai softmax dan cross-entropy, sehingga peluang semua kelas berjumlah satu. Kalau sebuah cerita
bisa sekaligus lucu, sedih, dan misterius, setiap label mendapat sigmoid sendiri dengan binary cross-entropy, dan setiap label
diputuskan dengan ambangnya sendiri. Memakai softmax untuk masalah multi-label memaksa label-label bersaing, padahal mereka tidak
saling meniadakan.

\section{Bahasa bersama modalitas lain}

NLP multimodal menggabungkan teks dengan gambar, suara, video, atau data sensor. Pada \istilah{visual question answering}
\citep{antol2015}, model menerima sebuah gambar dan pertanyaan tentangnya, misalnya ``berapa banyak orang di foto ini?'', dan harus
menjawab dalam bahasa. Pembuatan keterangan gambar berjalan ke arah sebaliknya. Kuncinya adalah ruang embedding bersama, tempat teks
dan gambar yang bermakna sama berdekatan, gagasan yang sama dengan CLIP di \cref{ch:hayati}. Kuliah juga menyebut penerapan seperti
kendaraan otonom yang memadukan bahasa dengan lidar dan kamera, dan komputasi afektif yang membaca emosi dari wajah, suara, dan kata.

\section{Biaya dan bias}

Kuliah tidak berhenti pada kemampuan. Model besar membutuhkan energi yang besar untuk dilatih, dan kefasihannya bisa menutupi
kenyataan bahwa ia tidak selalu memahami apa yang ditulisnya, argumen yang diangkat \citet{bender2021}. Bias sosial juga terbawa
dari data. Contoh dari kuliah adalah terjemahan mesin dari bahasa dengan kata ganti netral gender, seperti bahasa Turki atau bahasa
Indonesia, ke bahasa Inggris: ``dia seorang dokter'' dan ``dia seorang perawat'' bisa diterjemahkan menjadi \emph{he} dan \emph{she}
menurut stereotip, padahal kalimat aslinya tidak menyebut gender sama sekali. Mesin pencari gambar untuk kata ``CEO'' dan
``perawat'' memperlihatkan pola yang sama. Bias embedding kata yang dibahas \citet{bolukbasi2016} adalah akar dari masalah ini,
dan persoalan yang lebih luas dibahas di \cref{ch:manusia}.

\section{Perkembangan riset}

Banyak perkembangan model bahasa besar sudah dibahas di \cref{ch:terbaru}: hukum penskalaan, penyelarasan dengan preferensi manusia,
model yang bernalar, dan agen. Untuk NLP sebagai bidang, satu pergeseran besar adalah bahwa banyak tugas yang dulu punya model
khusus, seperti klasifikasi sentimen, ekstraksi entitas, atau terjemahan, kini dikerjakan oleh satu model umum yang diberi instruksi
dalam bahasa biasa. Tugas-tugas kuliah ini, dari kantong kata sampai BERT yang disetel, menjadi cara terbaik untuk memahami apa yang
sebenarnya terjadi di dalam model umum itu, dan untuk menilai kapan model kecil yang khusus masih lebih tepat, lebih murah, dan lebih
mudah diperiksa.

\sumberbab{Bab ini mengikuti kuliah dan latihan Natural Language Processing (pra-pemrosesan dan kantong kata, representasi kata,
PMI dan GloVe, jaringan saraf dan softmax, model bahasa $n$-gram, embedding kata dan kalimat yang dipelajari, attention dan
transformer, NLP multimodal, model bahasa besar). Buku rujukan yang disebut kuliah antara lain buku Jurafsky dan Martin tentang
pemrosesan ucapan dan bahasa.}
